% Lösungen Einheit 4 von 4 – Die Rate als Funktion (zusammenbau v0.1)
\einheitenkopf{Einheit 4 von 4 · Die Rate als Funktion}

\erg{19}{a) Bestand (Höhe, Anzahl, Menge) – größte Rate bei den Nullstellen der zweiten Ableitung: die Zuflussrate ist $f'$, ihr Maximum liegt bei $f'' = 0$ \quad b) $r'(t) = 0$ ⇔ $t = 6$ (der e-Faktor hat keine Nullstelle); nach 6 Stunden ist die Zuflussrate am größten \quad c) $f'(t) = 0$ ⇔ $t = 60$: am 60. Tag ist die Rate am größten (nicht $f(60)$ angeben) \quad d) $f'(x) = 1{,}5x^2 - 3$, $f''(x) = 3x = 0$ ⇔ $x = 0$; kleinste Steigung $f'(0) = -3$ (nicht das Minimum von $f$ suchen) \quad e) $w''(t) = -6t + 24 = 0$ ⇔ $t = 4$; $w'(4) = 48$ cm je Woche (nicht $t = 4$ als Ergebnis angeben) \quad f) $r'(t) = 50t \cdot (2 - 0{,}8t) \cdot \mathrm{e}^{-0{,}8t} = 0$ ⇔ $t = 0$ (Rand, Minimum) oder $t = 2{,}5$; $r(2{,}5) \approx 42{,}3$ Gäste je Stunde um 9:45 Uhr ($2{,}5$ h sind 2 h 30 min, nicht 2 h 50 min) \quad g) $z'(x) = 0{,}5 \cdot (6 - x) \cdot (6 - 3x) = 0$ ⇔ $x = 2$ oder $x = 6$; Vergleich: $z(0) = 0$, $z(2) = 16$, $z(6) = 0$ – größte Rate $16$ m³ je Stunde nach 2 Stunden, kleinste Rate $0$ zu Beginn und am Ende (Rand) \quad h) $h'(t) = 23{,}4$ ⇔ $t^2 - 20t + 36 = 0$ ⇔ $t = 2$ oder $t = 18$; $t = 18$ liegt außerhalb der Phase; $h'(0) = 45$: die Geschwindigkeit ist von $t = 0$ bis $t = 2$ mindestens $23{,}4$ m je min, der Zeitraum ist $2$ Minuten lang \quad i) $r(t) = 1{,}5t \cdot (10 - t)$ hat bei $t = 10$ eine Nullstelle mit Vorzeichenwechsel; für $t > 10$ ist $r(t) < 0$ – eine negative Anzahl von Bestellungen je Stunde gibt es nicht \quad j) Differenz der Raten $d(t) = a'(t) - b'(t) = 36t - 3t^2 - 6t = 30t - 3t^2$; $d'(t) = 30 - 6t = 0$ ⇔ $t = 5$; $d(5) = 75$ cm je Jahr; Ränder: $d(0) = 0$, $d(12) = -72$ – der Unterschied ist nach 5 Jahren mit $75$ cm je Jahr am größten (nicht die Differenz der Höhen untersuchen)}
\erg{20}{a) Fehler: Ben sucht das Maximum des Bestands $h$, nicht der Rate $h'$. Richtig: $h''(t) = -3t + 12 = 0$ ⇔ $t = 4$; $h'(4) = 24$ cm je Woche ist die größte Wachstumsgeschwindigkeit. \quad b) Die Änderungsrate ist $f'$; ihr Maximum liegt dort, wo $f''$ von plus nach minus wechselt – das ist eine Wendestelle von $f$, an der der Graph von links- in rechtsgekrümmt übergeht und am steilsten ist. \quad c) $c''(t) = 4 \cdot (0{,}25t - 1) \cdot \mathrm{e}^{-0{,}5t} = 0$ ⇔ $t = 4$; $c'(4) = -4 \cdot \mathrm{e}^{-2} \approx -0{,}54$ – vier Stunden nach der Einnahme sinkt die Konzentration mit etwa $0{,}54$ mg je Liter und Stunde am schnellsten \quad d) Hochpunkt des Graphen von $f'$: $t = 4$, Rate $= 8$ cm je Stunde}
